\documentclass[12pt]{article}
\usepackage[utf8]{inputenc}
\usepackage[T1]{fontenc}
\usepackage{hyperref}
\title{Kumbewaha Language Resources for K-12 Curriculum Development}
\author{Compiled by Claude Sonnet 5.5 (claude-sonnet-5-5) for Andr\'as Kornai}
\date{}
\begin{document}
\maketitle

\section{Introduction}
Kumbewaha (ISO 639-3 xks; Glottocode kumb1274) is an Austronesian language spoken on Buton Island, off the southeast coast of Sulawesi, Indonesia. Sources list about 3,400 speakers (2004). Joshua Project reports that all adults speak it but only about 40\% of children do. Little material exists for it. Search found no school materials, no published primers and no native-language curricular resources. A university news report describes Kumbewaha as an innovative dialect of Ciacia (Cia-Cia) that has been strongly influenced by the Lasalimu and Kamaru languages.

\section{Primary Educational Resources}
No primers, textbooks or school programmes for Kumbewaha were found in search results.

\section{Linguistic and Reference Materials}
\begin{itemize}
\item Wikipedia, ``Kumbewaha language'': general overview and classification. \url{https://en.wikipedia.org/wiki/Kumbewaha_language}
\item Austronesian Basic Vocabulary Database, Kumbewaha entry: a word list, placing the language in the Buton branch of the Muna-Buton group. \url{https://abvd.eva.mpg.de/austronesian/language.php?id=1521}
\item UGM news report on doctoral research on the Ciacia language, which mentions Kumbewaha as a Ciacia dialect (in Indonesian). \url{https://ugm.ac.id/id/berita/8742-raih-doktor-usai-teliti-bahasa-ciacia/}
\end{itemize}

\section{Religious and Community Resources}
\begin{itemize}
\item Joshua Project, Kumbewaha in Indonesia: demographic and location information (northeast coastal road between Pasarwajo and Lasalimu, Buton). It lists no known written form. \url{https://joshuaproject.net/people_groups/12844/ID}
\end{itemize}

\section{Digital and Audio Resources}
No audio, video, corpus or software resources specific to Kumbewaha were found in search results.

\section{Curriculum Development Context}
Kumbewaha has no known written standard and no known schooling use. A curriculum would first need community-led orthography work and documentation. The word lists and the closely related Ciacia language could serve as starting points, but the evidence found is thin.
\end{document}